\section{Rational characters from finite groups}\label{sec:characters}

Fix a prime $p$ and a positive height $l$, and put
$\one=L_{T(l)}S$.  Starting from a finite group $P$ with
$(U_P(\one))_{\Q}\ne0$, we will construct a pro-$p$ group
$\Gamma\twoheadrightarrow P$ whose finite sections require at most $l$
generators.  The same construction will describe rational cochain summands
by evaluation at continuous homomorphisms from $\Gamma$ to finite groups.
These evaluation coordinates will allow integration to impose the arithmetic conditions used in
\Cref{sec:profinite}.  The construction uses finite-groupoid ambidexterity
and elementary abelian sharpness, without requiring $\Gamma$ to be abelian.

\subsection{The rational tensor category and occurrence}

Let $\mathcal C$ be the homotopy category of dualizable objects in
$\Sp_{T(l)}$, with its local tensor product.  Write $\mathcal C_{\Q}$ for
the idempotent completion of the category with the same objects and
\[
 \Hom_{\mathcal C_{\Q}}(X,Y)
 =\Hom_{\mathcal C}(X,Y)\otimes\Q
 \qquad (X,Y\in\mathcal C).
\]
The tensor product extends to this category; its algebras are algebra
objects in this ordinary tensor category.  All algebra and module
identities below are inherited from the local category and then
rationalized.  In particular, this operation is distinct from ordinary
rationalization of spectra, which we denote by $(-)_{\Q}$.

For a finite groupoid $Y$, put $A_Y=\one^Y$, and abbreviate $A_{BF}$ to
$A_F$.  Finite-groupoid ambidexterity makes $A_Y$ dualizable and gives
natural algebra isomorphisms
\begin{equation}\label{eq:cochain-products}
 A_{Y\times Z}\simeq A_Y\otimes A_Z,
 \qquad
 A_{Y\sqcup Z}\simeq A_Y\times A_Z.
\end{equation}
Here and throughout this section tensors are local tensors before passing
to $\mathcal C_{\Q}$.  These facts follow from
\cite[Theorem~5.3.1 and Propositions~3.3.1, 3.3.7]{CSY}.
For a finite covering $f:Y\to Z$, the transfer $f_!:A_Y\to A_Z$ satisfies
base change and the projection formula.

Write $\mathbb A(F)$ for the Burnside ring of a finite group $F$.  Sending
an $F$-set $S$ to the transfer of $1$ along $S//F\to BF$ gives a ring map
\[
 \mathbb A(F)\longrightarrow\pi_0 A_F.
\]
The rational mark isomorphism
$\mathbb A(F)\otimes\Q\cong\prod_{(I\le F)}\Q$
has primitive idempotents $e_I^F$, characterized by the marks
$|(e_I^F)^J|=1$ if $J$ is conjugate to $I$ and $0$ otherwise.
Their images split $A_F$ into commutative algebra summands in
$\mathcal C_{\Q}$.  A finite group $L$ \emph{occurs} if
\[
 D_L:=e_L^L A_L\ne0.
\]
The unit of $D_L$ is the image of $e_L^L$ in
$\pi_0 A_L\otimes\Q$.  Since its identity is multiplication by this
class, $D_L\ne0$ precisely when this unit class is nonzero.

\begin{lemma}\label{lem:occurrence-fixed-points}
For every finite group $L$ and ordinary commutative ring spectrum $R$,
there is an equivalence of rational commutative ring spectra
\begin{equation}\label{eq:general-full-mark-comparison}
 e_L^L(R^{BL})_{\Q}\simeq\bigl(U_L(R)\bigr)_{\Q}.
\end{equation}
In particular,
\[
 D_L\ne0\quad\Longleftrightarrow\quad
 \bigl(U_L(\one)\bigr)_{\Q}\ne0.
\]
\end{lemma}
\begin{proof}
For any genuine $L$-spectrum $X$, categorical fixed points commute with
ordinary rationalization: their homotopy groups are represented by
suspensions of the compact genuine sphere.  In the rational genuine
category, $e_L^L X_{\Q}$ restricts to zero at every proper subgroup, since
$e_L^L$ does.  The isotropy separation sequence for the family of proper
subgroups consequently identifies
\[
 (e_L^L X_{\Q})^L\simeq
 \Phi^L(e_L^L X_{\Q})\simeq(\Phi^L X)_{\Q}.
\]
Apply this to the already formed cofree spectrum $X=\Bor_L(R)$, whose
categorical fixed points are $R^{BL}$.  The comparison respects
multiplication and identifies the full idempotent unit with the unit
of $(U_L(R))_{\Q}$, proving
\eqref{eq:general-full-mark-comparison}.  A unital ring spectrum is zero
exactly when its unit is zero.  Taking $R=\one$ and using the preceding
unit criterion for $D_L$ proves the last assertion.
\end{proof}

\begin{proposition}\label{prop:occurrence-sections}
Every occurring finite group is a $p$-group.  The class of occurring
groups is closed under subgroups and quotients.  If $A$ is any section
of an occurring group, its minimum number of generators satisfies
$d(A)\le l$.
\end{proposition}
\begin{proof}
We first record how a subgroup summand detects occurrence.  If $I\le F$,
the mark formula gives
\begin{equation}\label{eq:induced-top-mark}
 \operatorname{ind}_I^F(e_I^I)
   =|N_F(I):I|\,e_I^F.
\end{equation}
Indeed a subgroup has a nonzero mark on the left only when it is
conjugate to $I$, and the $I$-mark is $|N_F(I):I|$.
Covering transfer carries the full $I$-unit to the right side of
\eqref{eq:induced-top-mark}.  Thus $e_I^F A_F\ne0$ implies $D_I\ne0$.

Let $S$ be a Sylow $p$-subgroup of $F$.  The class $[F/S]\in\pi_0 A_F$
is a unit.  To see this, regard it as an endomorphism of the constant
unit in $\Fun(BF,\Sp_{T(l)})$.  Its value at the point is multiplication
by $[F:S]$, a $p$-local unit, and equivalences of diagrams are detected
at the point.  If $F$ is not a $p$-group, the full $F$-mark of $[F/S]$
is zero.  Hence $e_F^F[F/S]=0$, and invertibility forces $D_F=0$.

For a quotient $q:L\twoheadrightarrow Q$, pullback of $e_Q^Q$ has full
$L$-mark $1$.  Its image therefore acts as the identity on $D_L$.
If the full $Q$-unit were zero, its pullback would be zero, contradicting
occurrence of $L$.  This proves quotient closure.

For subgroup closure, let $I\le L$.  Before any rationalization, there
is a unit-preserving map of spectra
\begin{equation}\label{eq:subgroup-unit-map}
 U_I(\one)\longrightarrow U_L(\one).
\end{equation}
Here is its construction.  The cofree adjunction gives a map
\[
 N_I^L\Bor_I(\one)\longrightarrow\Bor_L(\one)
\]
whose adjoint on underlying spectra with $L$-action is indexed
multiplication
$\bigwedge_{L/I}\one\to\one$.
This uses the ordinary commutative ring spectrum underlying the local
unit; the displayed smash is the ordinary indexed smash.  The
multiplication is $L$-equivariant and carries the indexed unit to the
unit.  Apply geometric fixed points and the natural norm formula
$\Phi^L N_I^L\simeq\Phi^I$
\cite[Proposition~2.57]{HHR}.  This gives
\eqref{eq:subgroup-unit-map}, with the stated unit preservation.  If the
source unit were zero rationally, its image would be zero rationally.
\Cref{lem:occurrence-fixed-points} now proves subgroup closure.

It remains to bound generators.  If an occurring elementary abelian
group $E=(C_p)^a$ had $a>l$, the elementary abelian sharpness theorem
of Arone, Dwyer, Lesh, and Mahowald, in the form
\cite[Theorems~2.2 and~4.1]{BHNNNS}, would give a finite genuine $E$-spectrum $X$
with underlying type $a$ and full geometric fixed points of type $0$.
Its underlying spectrum is $T(l)$-acyclic.  Finite duality and the
cofree projection formula give
\[
 X\wedge\Bor_E(\one)
   \simeq\Bor_E\bigl((\operatorname{res}X)\wedge\one\bigr)=0.
\]
Here $\operatorname{res}X$ retains its underlying $E$-action.  Its smash
with $\one$ is zero because this can be tested on the underlying
spectrum, and smashing a finite spectrum with the local unit computes
its localization.  Taking geometric fixed points and
rationalizing yields
\[
 (\Phi^E X)_{\Q}\wedge(U_E(\one))_{\Q}=0.
\]
The first factor is a nonzero finite rational spectrum, so it detects
nonzero rational spectra under smash.  This contradicts occurrence of
$E$.  Finally, every section $A$ of an occurring group occurs, and so
does its Frattini quotient $A/\Phi(A)\cong(C_p)^{d(A)}$.  The elementary
abelian bound and the Burnside basis theorem give $d(A)\le l$.
\end{proof}

\subsection{A maximizing tower}

For a profinite group $\Gamma$ and a finite groupoid $Y$, let
\[
 \mathcal X_\Gamma(Y)=\Fun_{\cts}(B\Gamma,Y).
\]
Continuity means that the homomorphism into the automorphism group of
the chosen object has open kernel.  Thus
$\mathcal X_\Gamma(BF)=\Hom_{\cts}(\Gamma,F)//F$.
A finite section of $\Gamma$ means a finite continuous quotient of a
closed subgroup.

\begin{theorem}[Pro-character construction]\label{thm:pro-characters}
Suppose $P$ is a finite group with $(U_P(\one))_{\Q}\ne0$.
There is a tower of occurring finite $p$-groups
\[
 \cdots\twoheadrightarrow L_2\twoheadrightarrow L_1
 \twoheadrightarrow L_0=P,
 \qquad \Gamma=\varprojlim_i L_i,
\]
with the following properties.
\begin{enumerate}[label=\textup{(\roman*)}]
\item Every finite section $A$ of $\Gamma$ has $d(A)\le l$.
Every closed subgroup of $\Gamma$ is topologically generated by at most
$l$ elements; in particular $\Gamma$ is finitely generated and
$\Gamma\twoheadrightarrow P$.
\item For each finite collection of finite groupoids $Y$, at every
sufficiently advanced stage $L=L_i$ inflation is an equivalence
\begin{equation}\label{eq:mapping-groupoid-stabilization}
 \Fun(BL,Y)\xrightarrow{\ \sim\ }\mathcal X_\Gamma(Y).
\end{equation}
At those stages the evaluation maps give an isomorphism of
$D_L$-algebras in $\mathcal C_{\Q}$,
\begin{equation}\label{eq:character-splitting}
 D_L\otimes A_Y\xrightarrow{\ \sim\ }
 \prod_{[\alpha]\in\pi_0\mathcal X_\Gamma(Y)}D_L.
\end{equation}
These isomorphisms are natural for pullback along maps of the selected
groupoids.  In particular, a common stage can be chosen for a finite
diagram of finite groupoids, including both its bases and its total
groupoids.
\end{enumerate}
\end{theorem}
\begin{proof}
Enumerate the isomorphism classes of finite groups as $F_0,F_1,\ldots$.
Given an occurring group $L_i$, choose an occurring extension
$L_{i+1}\twoheadrightarrow L_i$ maximizing
\[
 \bigl|\Hom(L_{i+1},F_i)/F_i\bigr|
\]
among all occurring extensions of $L_i$.  The set of candidates is
nonempty, since $L_i$ is one.  By
\Cref{prop:occurrence-sections}, every candidate has at most $l$
generators, so the displayed integer is at most $|F_i|^l$ and a maximum
is attained.  Precomposition with a surjection injects conjugacy
classes of homomorphisms.  Once a target $F_i$ has been maximized, every
later occurring extension is still an extension of $L_i$ and cannot
increase the maximum.  Consequently, for every sufficiently advanced
stage $L$ and every occurring extension $M\twoheadrightarrow L$,
\begin{equation}\label{eq:maximal-hom-classes}
 \Hom(L,F_i)/F_i\xrightarrow{\ \sim\ }\Hom(M,F_i)/F_i.
\end{equation}

Every continuous homomorphism $\Gamma\to F$ factors through some
$L_i$, since the stage kernels form a neighborhood basis of the
identity.  After $F$ has been maximized, the preceding bijections imply
that every such homomorphism factors through each later stage.  This
is actual factorization: a map conjugate to one factoring through a
quotient kills the same quotient kernel.  Surjectivity of
$\Gamma\to L_i$ also identifies all conjugating elements.  This proves
\eqref{eq:mapping-groupoid-stabilization} first for $Y=BF$, and then
for finite disjoint unions of such groupoids.

If $H$ is closed in $\Gamma$ and $H/K$ is finite, choose an open normal
subgroup $N\trianglelefteq_o\Gamma$ with $H\cap N\subseteq K$.
Then $H/K$ is a quotient of the subgroup $H/(H\cap N)$ of $\Gamma/N$.
The latter is a quotient of a tower stage, so
\Cref{prop:occurrence-sections} gives $d(H/K)\le l$.
For each open normal $K$ of $H$, the tuples in $H^l$ whose images
generate $H/K$ form a nonempty clopen subset.  These subsets have the
finite intersection property, by applying the same bound to the
quotient by an intersection of finitely many such $K$.  Compactness
gives a tuple generating every finite quotient, hence topologically
generating $H$.  This proves \textup{(i)}.

We prove the splitting for $Y=BF$, at a stage $L$ satisfying
\eqref{eq:maximal-hom-classes} for $F$.  We first show that every surviving
subgroup summand comes from a graph, and then identify each graph summand
by restriction.  Put $G=L\times F$ and let
$\pi:G\to L$ be the first projection.  By
\eqref{eq:cochain-products},
$D_L\otimes A_F$ is the summand $\pi^*(e_L^L)A_G$.
The marks of $\pi^*(e_L^L)$ are supported exactly on subgroups
$M\le G$ projecting onto $L$.  If $e_M^G A_G\ne0$, then $M$ occurs by
\eqref{eq:induced-top-mark}.  It is therefore an occurring extension of
$L$.  Apply \eqref{eq:maximal-hom-classes} to its second projection
$M\to F$: this map is conjugate to $\alpha\circ\pi$ for some
$\alpha:L\to F$.  Absorbing the conjugation into $\alpha$ shows that
$M$ is contained in its graph, and surjectivity onto $L$ gives
\[
 M=M_\alpha:=\{(x,\alpha(x)):x\in L\}.
\]
Graphs are conjugate in $L\times F$ exactly when the corresponding
homomorphisms are conjugate in $F$: conjugation in the first factor is
already absorbed by conjugation by its image under $\alpha$.

We next identify each graph summand, keeping track of the map.
Let $B_\alpha=e_{M_\alpha}^G A_G$ and write
$j_\alpha:L\xrightarrow{\sim}M_\alpha\hookrightarrow G$.
Restriction along $j_\alpha$, followed by the full projection, gives
\[
 r_\alpha:B_\alpha\longrightarrow D_L.
\]
It has a right inverse $s_\alpha$, obtained by pulling back along
$\pi$ and projecting to $B_\alpha$, because
$\pi j_\alpha=\id$ and
$j_\alpha^*(e_{M_\alpha}^G)=e_L^L$.
For the other inverse, let $t_\alpha:D_L\to B_\alpha$ be covering
transfer along $j_\alpha$, followed by projection.  The projection
formula gives
\[
 t_\alpha r_\alpha
  = |(G/M_\alpha)^{M_\alpha}|\,\id_{B_\alpha}
  = |C_F(\alpha(L))|\,\id_{B_\alpha}.
\]
Indeed the full graph mark of $[G/M_\alpha]$ is
$|N_G(M_\alpha):M_\alpha|=|C_F(\alpha(L))|$.
This positive integer is invertible in $\mathcal C_{\Q}$, so
$|C_F(\alpha(L))|^{-1}t_\alpha$ is a left inverse.
Thus $r_\alpha$ is an isomorphism, with inverse $s_\alpha$.
In particular every graph summand survives, since $D_L\ne0$.

Combining the graph summands proves
\eqref{eq:character-splitting} for $BF$.  Its coordinate at
$\alpha:L\to F$ is explicitly the diagonal restriction
\begin{equation}\label{eq:character-evaluation}
 D_L\otimes A_F\longrightarrow D_L,
 \qquad d\otimes a\longmapsto d\,\bigl(e_L^L\alpha^*a\bigr).
\end{equation}
It is a map of $D_L$-algebras.  Composition of homomorphisms composes
these restrictions, giving naturality.  Applying the disjoint-union
identity in \eqref{eq:cochain-products} proves the assertion for every
finite groupoid.  Only finitely many component automorphism groups
occur in any finite diagram, so they can all be maximized before
choosing the stage.  This proves \textup{(ii)}.
\end{proof}

Naturality also controls homotopy fibers. For a map $Y\to Z$ of finite
groupoids, choose a stage satisfying
\eqref{eq:mapping-groupoid-stabilization} for both $Y$ and $Z$.
The resulting square of mapping groupoids is a square of equivalences
on its horizontal arrows.  It therefore identifies homotopy fibers
over any chosen base label.  For a groupoid bundle over $BQ$, the
targets to be stabilized include $Q$ and the automorphism groups of
the components of its \emph{total} groupoid.  Stabilizing only the
quotient $Q$ defining the action would not ensure that all continuous
lifts of the base label are counted.

\begin{lemma}[Evaluation of covering transfers]\label{lem:cover-evaluation}
Let $f:Y\to Z$ be a finite covering of finite groupoids.  Choose a
common stage $L$ as in \Cref{thm:pro-characters}, and a label
$\alpha:BL\to Z$.  If the pullback of $f$ is the finite $L$-set $S$,
then the evaluation of $f_!(1)\in\pi_0 A_Z$ at $\alpha$ is
$|S^L|\,1_{D_L}$.
\end{lemma}
\begin{proof}
Base change identifies $\alpha^*f_!(1)$ with the Burnside class $[S]$
in $\pi_0 A_L$.  Multiplication by $e_L^L$ selects its full mark
$|S^L|$.  This is exactly the evaluation map
\eqref{eq:character-evaluation}.
\end{proof}
